\clearpage\phantomsection\label{preuves-debut}\pdfbookmark[0]{Démonstrations de référence}{bm-preuves}
\editionPage{preuves-1}{Démonstrations}{Multiples et parité}{Lire les hypothèses, les transformations et la conclusion.}
\begin{blocEdition}{preuves-1}
\begin{demonstration}[1. Somme de deux multiples de 7]
Si $a=7m$ et $b=7n$ avec $m,n\in\mathbb Z$, alors $a+b=7(m+n)$. Comme $m+n$ est entier, la somme est un multiple de 7. Le même calcul fonctionne avec tout diviseur fixé.
\end{demonstration}
\begin{demonstration}[2. Carré d’un impair]
Pour $n=2k+1$, $n^2=4k^2+4k+1=2(2k^2+2k)+1$. Le nombre $2k^2+2k$ est entier, donc $n^2$ est impair. La contraposée donne : si $n^2$ est pair, alors $n$ est pair.
\end{demonstration}
\begin{demonstration}[3. Pourquoi un tiers n’est pas décimal]
Supposons $1/3=N/10^k$ pour un entier $N$ et un entier $k\ge0$. Alors $10^k=3N$ serait divisible par 3. Or $10^k$ a une écriture « 1 suivi de zéros », dont la somme des chiffres vaut 1 ; il n’est pas divisible par 3. La contradiction prouve que $1/3$ n’est pas décimal.
\end{demonstration}

\end{blocEdition}
\editionPage{preuves-2}{Démonstrations}{Racines et irrationalité}{Lire les hypothèses, les transformations et la conclusion.}
\begin{blocEdition}{preuves-2}
\begin{demonstration}[4. Irrationalité de la racine de 2]
Supposons $\sqrt2=p/q$ avec $p,q$ entiers positifs et la fraction irréductible. En élevant au carré, $p^2=2q^2$, donc $p^2$ est pair. Par la propriété précédente, $p$ est pair : $p=2r$. Alors $4r^2=2q^2$, donc $q^2=2r^2$, et $q$ est pair aussi. Ainsi 2 divise $p$ et $q$, ce qui contredit l’irréductibilité. Le nombre $\sqrt2$ est irrationnel.
\end{demonstration}
\begin{demonstration}[5. Racine d’un produit]
Pour $a,b\ge0$, les nombres $\sqrt{ab}$ et $\sqrt a\sqrt b$ sont positifs ou nuls. De plus $(\sqrt a\sqrt b)^2=ab$. Par unicité de la racine carrée positive de $ab$, on obtient $\sqrt{ab}=\sqrt a\sqrt b$. La condition de positivité distingue la racine positive de son opposé.
\end{demonstration}

\end{blocEdition}
\editionPage{preuves-3}{Démonstrations}{Vecteurs et équations de droite}{Lire les hypothèses, les transformations et la conclusion.}
\begin{blocEdition}{preuves-3}
\begin{demonstration}[6. Déterminant et colinéarité]
Si $\vec v=k\vec u$, alors pour $\vec u=(a;b)$ et $\vec v=(c;d)$, $ad-bc=a(kb)-b(ka)=0$. Réciproquement, supposons $\vec u\ne\vec0$ et $ad-bc=0$. Si $a\ne0$, poser $k=c/a$ ; l’égalité donne $d=bc/a=kb$, donc $\vec v=k\vec u$. Si $a=0$, alors $b\ne0$ et $-bc=0$ impose $c=0$ ; poser $k=d/b$ donne encore $\vec v=k\vec u$. Le cas d’un vecteur nul est colinéaire par convention.
\end{demonstration}
\begin{demonstration}[7. Forme cartésienne d’une droite]
Une droite passant par $A(x_0;y_0)$ de directeur non nul $(u;v)$ contient exactement les $M(x;y)$ tels que $u(y-y_0)-v(x-x_0)=0$. Cela s’écrit $-vx+uy+vx_0-uy_0=0$, de la forme $ax+by+c=0$ avec $(a;b)\ne(0;0)$. Réciproquement, pour une telle équation, choisir un point solution $A$ puis utiliser le directeur $(-b;a)$ ; le même calcul décrit une droite. Un point solution existe : si $b\ne0$, prendre $(0;-c/b)$ ; sinon $(-c/a;0)$.
\end{demonstration}

\end{blocEdition}
\editionPage{preuves-4}{Démonstrations}{Variations et comparaison}{Lire les hypothèses, les transformations et la conclusion.}
\begin{blocEdition}{preuves-4}
\begin{demonstration}[8. Variations d’une fonction affine]
Pour $f(x)=mx+p$ et $a<b$, $f(b)-f(a)=m(b-a)$. Le facteur $b-a$ est positif. La différence a donc le signe de $m$ : positive pour $m>0$, négative pour $m<0$, nulle pour $m=0$. Cela prouve les trois affirmations de variation.
\end{demonstration}
\begin{demonstration}[9. Position du carré par rapport à l’identité]
Pour $x\ge0$, $x^2-x=x(x-1)$. Sur $[0;1]$, le produit est négatif ou nul ; sur $[1;+\infty[$, il est positif ou nul. Les deux courbes sont donc dans l’ordre $x^2\le x$ entre 0 et 1, puis dans l’ordre inverse après 1. Les égalités sont obtenues en 0 et 1.
\end{demonstration}

\end{blocEdition}
\editionPage{preuves-5}{Démonstrations}{Variations des fonctions de référence}{Lire les hypothèses, les transformations et la conclusion.}
\begin{blocEdition}{preuves-5}
\begin{demonstration}[10. La fonction carré]
Pour $0\le a<b$, $b^2-a^2=(b-a)(b+a)>0$ : le carré est strictement croissant sur $[0;+\infty[$. Pour $a<b\le0$, le premier facteur est positif et $a+b<0$, donc $b^2-a^2<0$ : le carré est strictement décroissant sur $]-\infty;0]$. Il faut traiter séparément les deux intervalles.
\end{demonstration}
\begin{demonstration}[11. La fonction inverse]
Pour $a<b$ tous deux strictement positifs, ou tous deux strictement négatifs, $ab>0$. On a $1/b-1/a=(a-b)/(ab)<0$. La fonction inverse est donc strictement décroissante sur chacun des deux intervalles de son domaine. L’argument ne vaut pas pour un couple qui traverse 0, car le signe de $ab$ change.
\end{demonstration}

\end{blocEdition}
\editionPage{preuves-6}{Démonstrations}{Coordonnées et statistiques}{Lire les hypothèses, les transformations et la conclusion.}
\begin{blocEdition}{preuves-6}
\begin{demonstration}[12. Coordonnées du milieu]
Pour un milieu $I$, $\vect{AI}=\vect{IB}$. Les abscisses donnent $x_I-x_A=x_B-x_I$, soit $2x_I=x_A+x_B$. De même, $2y_I=y_A+y_B$. Réciproquement, ces coordonnées donnent l’égalité des deux vecteurs et caractérisent donc le milieu. La formule est ainsi démontrée dans les deux sens.
\end{demonstration}
\begin{demonstration}[13. Deux écritures de la dispersion]
En développant et en prenant la moyenne des carrés des écarts, on obtient $\overline{(x-\bar x)^2}=\overline{x^2}-2\bar x\,\overline{x}+\bar x^2$. Comme $\overline{x}=\bar x$, cela donne $\overline{x^2}-\bar x^2$. Les pondérations éventuelles restent les mêmes dans toutes les moyennes. Cette identité justifie la formule utilisée dans les programmes.
\end{demonstration}

\end{blocEdition}
